% Generated draft: requires mathematical and source review.
\documentclass[11pt]{article}
\usepackage[margin=25mm]{geometry}
\usepackage{fontspec}
\setmainfont{Latin Modern Roman}
\usepackage{amsmath,amssymb,mathtools,unicode-math}
\setmathfont{Latin Modern Math}
\usepackage{enumitem,needspace,xurl,hyperref}
\setlength{\parindent}{0pt}
\setlength{\parskip}{5pt}
\setlength{\emergencystretch}{3em}
\newcommand{\N}{\mathbb{N}}
\newcommand{\Z}{\mathbb{Z}}
\newcommand{\Q}{\mathbb{Q}}
\newcommand{\R}{\mathbb{R}}
\newcommand{\C}{\mathbb{C}}
\newcommand{\F}{\mathbb{F}}
\newcommand{\A}{\mathbb{A}}
\newcommand{\PP}{\mathbb{P}}
\newcommand{\E}{\mathbb{E}}
\newcommand{\eps}{\varepsilon}
\newcommand{\dd}{\,\mathrm{d}}
\begin{document}
\section*{UnsolvedMath ID: 30000652}
\textbf{Source problem code:} catalogue entry
\subsection{Definitions and mathematical statement}
Let \(X\) be a smooth projective complex rational surface, so \(X\) is birational to \(\PP^2_\C\). Let \(\operatorname{Pic}(X)\) be its group of line bundles, and define
\[
\operatorname{Aut}^*(X)=\operatorname{im}\bigl(\operatorname{Aut}(X)\to\operatorname{Aut}(\operatorname{Pic}(X))\bigr)
\]
using a consistent pullback-action convention. Assume this image is infinite. Infinitude of the automorphism group alone is not substituted for infinitude of its Picard action.

The source asks whether there is a birational morphism \(\pi:X\to Y\) such that \(Y\) has an antipluricanonical curve and an infinite subgroup of \(\operatorname{Aut}^*(Y)\) lifts to \(X\). Under the usual smooth-surface convention, the first condition means
\[
\exists m\geq1\quad\exists 0\ne D\geq0,\qquad D\sim-mK_Y.
\]
No irreducibility, reducedness, or smoothness of the divisor is requested. The condition is on the target \(Y\), not necessarily on \(X\). The source does not explicitly state whether \(Y\) is smooth projective or permits singular targets; the latter would require a specified canonical-divisor and pluricanonical-sheaf convention.

To make the usual lifting requirement precise, set
\[
L_\pi=\{a\in\operatorname{Aut}(Y):\exists\widetilde a\in\operatorname{Aut}(X),\quad
\pi\circ\widetilde a=a\circ\pi\}.
\]
The intended condition can be expressed as infinitude of the image of \(L_\pi\) in \(\operatorname{Aut}^*(Y)\). This gives actual liftable automorphisms whose Picard actions form an infinite subgroup; it does not require a section from all of \(\operatorname{Aut}^*(Y)\) or lifting every automorphism of \(Y\). If the source intends an action of a specified abstract subgroup with chosen representatives, that stronger convention must be recovered.

Thus for every indicated \(X\) with infinite Picard-action image, the desired morphism and effective multiple of \(-K_Y\) may depend on \(X\), while the lifted Picard-action image must still be infinite. The supplied literature triage checked 2026-08-21 reports no authoritative complete resolution of the extracted question, without a separate theorem or construction.
\subsection{Short English statement}
Can a rational surface with infinitely many Picard-group automorphisms be contracted to a surface carrying an effective anticanonical multiple while retaining an infinite liftable Picard-action subgroup? The target and lifting conventions need confirmation.
\subsection{Sources}
\begin{itemize}
\item Local UnsolvedMath registry, permanent ID 30000652, catalogue entry, \emph{Automorphisms of Rational Surfaces with Antipluricanonical Curves}.
\item \emph{Affine Algebraic Geometry} (2007), Oberwolfach Reports, pp. 5--82, DOI \url{https://doi.org/10.4171/owr/2007/01}.
\item Supplied source-verification claim and literature triage checked 2026-08-21 cite only the report; no separate contraction or lifting theorem is supplied.
\end{itemize}
\end{document}
